\begin{proposition}[Submodule criterion] \label{proposition:submodule_criterion_for_a_subset_of_a_module_over_a_ring}
    Let $R,S$ be \CrefAndHyperrefIfExist{definition:ring}{not-necessarily commutative rings} and let $M$ be an \CrefAndHyperrefIfExist{definition:module_of_a_ring}{$R$-$S$-bimodule}. Let $N \subseteq M$ be a subset. Then $N$ is an \CrefAndHyperrefIfExist{definition:submodule_of_a_module_over_a_ring}{$R$-$S$-submodule of $M$} if and only if the following conditions hold:
    \begin{itemize}
        \item $N \neq \emptyset$,
        \item For all $n_1, n_2 \in N$, the sum $n_1 + n_2$ is an element of $N$,
        \item For all $r \in R$ and $n \in N$, the product $r n$ is an element of $N$,
        \item For all $s \in S$ and $n \in N$, the product $n s$ is an element of $N$.
    \end{itemize}
    In particular, 
    \begin{enumerate}
        \item For a left $R$-module $M$, a subset $N \subseteq M$ is a left $R$-submodule if and only if it is nonempty, closed under addition, and closed under left multiplication by $R$
        \item For a right $R$-module $M$, a subset $N \subseteq M$ is a right $R$-submodule if and only if it is nonempty, closed under addition, and closed under right multiplication by $R$
        \item For a two-sided $R$-module $M$, a subset $N \subseteq M$ is a two-sided $R$-submodule if and only if it is nonempty, closed under addition, and closed under both left and right multiplication by $R$.
    \end{enumerate}
\end{proposition}
\begin{proof}
    If $N$ is an $R$-$S$-submodule of $M$, then the listed conditions hold by definition. Conversely, suppose that the listed conditions hold. We show that $N$ is indeed an $R$-$S$-submodule of $M$.

    First, we verify that $N$ is a subgroup of the underlying abelian group of $M$. By \CrefAndHyperrefIfExist{proposition:subgroup_criterion_for_a_subset_of_a_group}{the subgroup criterion}, it suffices to show that $N$ is nonempty, closed under addition, and closed under additive inverses. The first two conditions are given. For closure under additive inverses, fix any $n \in N$. Since $(-1) \cdot n = -n$\CrefIfExists{lemma:minus_one_times_element_is_inverse_in_Z_module} and $N$ is closed under left multiplication by elements of $R$, we have $-n = (-1) \cdot n \in N$. Therefore, $N$ is a subgroup of $M$.

    Next, $N$ is closed under left multiplication by $R$ and right multiplication by $S$ by assumption. Hence, $N$ is an $R$-$S$-submodule of $M$ as desired.
\end{proof}